\documentclass[10pt, oneside]{article}   	% use "amsart" instead of "article" for AMSLaTeX format
\usepackage{geometry}                		% See geometry.pdf to learn the layout options. There are lots.
\geometry{letterpaper}                   		% ... or a4paper or a5paper or ... 
%\geometry{landscape}                		% Activate for rotated page geometry
%\usepackage[parfill]{parskip}    		% Activate to begin paragraphs with an empty line rather than an indent
\usepackage{graphicx}				% Use pdf, png, jpg, or eps§ with pdflatex; use eps in DVI mode
								% TeX will automatically convert eps --> pdf in pdflatex		
\usepackage{amssymb}
\usepackage{mhchem}
\usepackage{hyperref}

\title{Module 14, Part 1}
\author{Mark Peever\\ \texttt{mpeever@gmail.com}}
\date{April 10, 2026}

\begin{document}
\maketitle

\section*{Objectives}
\marginpar{0 minutes}
Refer to \href{https://drive.google.com/file/d/1Bxe0uxWzNG95BBQg0zQ9IEhdGUAHYCXS/view?usp=sharing}{Module 14 Notes}.\\

By the end of this class, the students should be able to\ldots
\begin{itemize}
\item describe what a rate is
\item describe what a reaction rate is
\item describe and explain the Rate Equation
\item describe what a catalyst is
\item list some contributing factors to reaction rates
\end{itemize}

\section*{Welcome \& Devotion}
\marginpar{5 minutes}
\begin{itemize}
\item have one student read \href{https://www.biblegateway.com/passage/?search=psalm\%20107\&version=LSB}{Psalm 107:17--22}
\end{itemize}


\section*{Module 13 Quiz}
\marginpar{20 minutes}


\section*{Introduction to Kinetics}
\marginpar{5 minutes}

\begin{itemize}
\item \emph{Kinetics}  is all about the rate of chemical reactions
\item describe what a \emph{rate} is
\item $R = \frac{\Delta [product]}{\Delta t} = - \frac{\Delta [reactant]}{\Delta t}$
\item describe how these are really differential equations 
\end{itemize}

\section*{Qualitative Kinetics}
\marginpar{10 minutes}

\begin{itemize}
\item reactions are faster at higher temperatures
\item reactions are faster with more surface area
\item reactions are faster with increased concentration
\item \emph{Why?}
\end{itemize}

\section*{Quantitative Kinetics}
\marginpar{10 minutes}

\begin{itemize}
\item $R = k(t)  [A]^{x}  [B]^{y}$
\item go through examples!
\end{itemize}

\section*{Catalysts}
\marginpar{10 minutes}
\begin{itemize}
\item catalysts improve reaction rates without themselves being reactants
\item catalysts lower activation energy
\item sometimes they're used up in one step and then produced again in another
\item sometimes they're not reactants at all, but bring reactants closer together, etc.
\item explain homogeneous vs. heterogeneous catalysts
\end{itemize}

\section*{Questions for me}
\marginpar{5 minutes}

\section*{Assignment}
\marginpar{5 minutes}
\begin{itemize}
\item Module 14 Review Problems: p. 488 \# 1--10 (not to be turned in)
\item Module 14 Practice Problems: pp. 489--490 \#1--10 due 2026-04-17
\end{itemize}



\end{document}  